%%% Hlavní soubor. Zde se definují základní parametry a odkazuje se na ostatní části. %%%

%% Verze pro jednostranný tisk:
% Okraje: levý 40mm, pravý 25mm, horní a dolní 25mm
% (ale pozor, LaTeX si sám přidává 1in)
\documentclass[12pt,a4paper]{report}
\setlength\textwidth{145mm}
\setlength\textheight{247mm}
\setlength\oddsidemargin{15mm}
\setlength\evensidemargin{15mm}
\setlength\topmargin{0mm}
\setlength\headsep{0mm}
\setlength\headheight{0mm}
% \openright zařídí, aby následující text začínal na pravé straně knihy
\let\openright=\clearpage

%% Pokud tiskneme oboustranně:
% \documentclass[12pt,a4paper,twoside,openright]{report}
% \setlength\textwidth{145mm}
% \setlength\textheight{247mm}
% \setlength\oddsidemargin{15mm}
% \setlength\evensidemargin{0mm}
% \setlength\topmargin{0mm}
% \setlength\headsep{0mm}
% \setlength\headheight{0mm}
% \let\openright=\cleardoublepage

%% Pokud používáte csLaTeX (doporučeno):
%\usepackage{czech}
%% Pokud nikoliv:
\usepackage[czech]{babel}
\usepackage[T1]{fontenc}

%% Použité kódování znaků: obvykle latin2, cp1250 nebo utf8:
\usepackage[utf8]{inputenc}

%% Ostatní balíčky
\usepackage{graphicx}
\usepackage{amsthm}

%% Balíček hyperref, kterým jdou vyrábět klikací odkazy v PDF,
%% ale hlavně ho používáme k uložení metadat do PDF (včetně obsahu).
%% POZOR, nezapomeňte vyplnit jméno práce a autora.
\usepackage[unicode]{hyperref}   % Musí být za všemi ostatními balíčky
\hypersetup{pdftitle=Vývoj hlasově ovládaných webových her pomocí CloudASR}
\hypersetup{pdfauthor=Jan Milota}

%%% Drobné úpravy stylu

% Tato makra přesvědčují mírně ošklivým trikem LaTeX, aby hlavičky kapitol
% sázel příčetněji a nevynechával nad nimi spoustu místa. Směle ignorujte.
\makeatletter
\def\@makechapterhead#1{
  {\parindent \z@ \raggedright \normalfont
   \Huge\bfseries \thechapter. #1
   \par\nobreak
   \vskip 20\p@
}}
\def\@makeschapterhead#1{
  {\parindent \z@ \raggedright \normalfont
   \Huge\bfseries #1
   \par\nobreak
   \vskip 20\p@
}}
\makeatother

% Toto makro definuje kapitolu, která není očíslovaná, ale je uvedena v obsahu.
\def\chapwithtoc#1{
\chapter*{#1}
\addcontentsline{toc}{chapter}{#1}
}

\begin{document}

% Trochu volnější nastavení dělení slov, než je default.
\lefthyphenmin=2
\righthyphenmin=2

%%% Titulní strana práce

\pagestyle{empty}
\begin{center}

\large

Univerzita Karlova v Praze

\medskip

Matematicko-fyzikální fakulta

\vfill

{\bf\Large BAKALÁŘSKÁ PRÁCE}

\vfill

\centerline{\mbox{\includegraphics[width=60mm]{img/logo.png}}}

\vfill
\vspace{5mm}

{\LARGE Jan Milota}

\vspace{15mm}

% Název práce přesně podle zadání
{\LARGE\bfseries Vývoj hlasově ovládaných webových her pomocí CloudASR}

\vfill

% Název katedry nebo ústavu, kde byla práce oficiálně zadána
% (dle Organizační struktury MFF UK)
Ústav formální a aplikované lingvistiky

\vfill

\begin{tabular}{rl}

Vedoucí bakalářské práce: & Mgr. Ing. Filip Jurčíček, Ph.D. \\
\noalign{\vspace{2mm}}
Studijní program: & Informatika \\
\noalign{\vspace{2mm}}
Studijní obor: & Obecná  Informatika \\
\end{tabular}

\vfill

% Zde doplňte rok
Praha 2015

\end{center}

\newpage

%%% Následuje vevázaný list -- kopie podepsaného "Zadání bakalářské práce".
%%% Toto zadání NENÍ součástí elektronické verze práce, nescanovat.

%%% Na tomto místě mohou být napsána případná poděkování (vedoucímu práce,
%%% konzultantovi, tomu, kdo zapůjčil software, literaturu apod.)

\openright

\noindent
Děkuji panu doktorovi Jurčíčkovi za profesionální vedení této práce a jeho neutuchající víru v mou schopnost samostatného vývoje.

\newpage

%%% Strana s čestným prohlášením k bakalářské práci

\vglue 0pt plus 1fill

\noindent
Prohlašuji, že jsem tuto bakalářskou práci vypracoval samostatně a výhradně
s~použitím citovaných pramenů, literatury a dalších odborných zdrojů.

\medskip\noindent
Beru na~vědomí, že se na moji práci vztahují práva a povinnosti vyplývající
ze zákona č. 121/2000 Sb., autorského zákona v~platném znění, zejména skutečnost,
že Univerzita Karlova v Praze má právo na~uzavření licenční smlouvy o~užití této
práce jako školního díla podle §60 odst. 1 autorského zákona.

\vspace{10mm}

\hbox{\hbox to 0.5\hsize{%
V ........ dne ............
\hss}\hbox to 0.5\hsize{%
Podpis autora
\hss}}

\vspace{20mm}
\newpage

%%% Povinná informační strana bakalářské práce

\vbox to 0.5\vsize{
\setlength\parindent{0mm}
\setlength\parskip{5mm}

Název práce:
Vývoj hlasově ovládaných webových her pomocí CloudASR
% přesně dle zadání

Autor:
Jan Milota

Katedra:  % Případně Ústav:
Ústav formální a aplikované lingvistiky
% dle Organizační struktury MFF UK

Vedoucí bakalářské práce:
Mgr. Ing. Filip Jurčíček, Ph.D., Ústav formální a aplikované lingvistiky
% dle Organizační struktury MFF UK, případně plný název pracoviště mimo MFF UK

Abstrakt:
% abstrakt v rozsahu 80-200 slov; nejedná se však o opis zadání bakalářské práce
Cílem práce je navrhnout a vyvinout software pro výuku jazyků hrou za použití webových technologií a~čerstvě vznikající CloudASR knihovny. Běžný uživatel provozuje 
interakci se svým prohlížečem skoro výhradně prostřednictvím myši a klávesnice. Díky softwaru, který tato práce reprezentuje, má nyní uživatel příležitost zabřednout do 
někdy ne úplně populární výuky jazyka i za pomoci svého hlasu. Což nabízí zmíněné výuce netušené možnosti, obzvláště stran uživatelské interaktivity. Důraz byl kladen 
na~uživatelskou přívětivost, grafickou fidelitu a~na~kompetitivní aspekt výuky využívajíc Facebookovou integraci a bodové hodnotící žebříčky.

Klíčová slova:
automatické rozpoznávání řeči, ASR, hry, web, HTML5, Javascript

\vss}\nobreak\vbox to 0.49\vsize{
\setlength\parindent{0mm}
\setlength\parskip{5mm}

Title:
% přesný překlad názvu práce v angličtině
Development of speech enabled web games using CloudASR

Author:
Jan Milota

Department:
Institute of Formal and Applied Linguistics
% dle Organizační struktury MFF UK v angličtině

Supervisor:
Mgr. Ing. Filip Jurčíček, Ph.D., Institute of Formal and Applied Linguistics
% dle Organizační struktury MFF UK, případně plný název pracoviště
% mimo MFF UK v angličtině

Abstract:
% abstrakt v rozsahu 80-200 slov v angličtině; nejedná se však o překlad
% zadání bakalářské práce
The main goal of this thesis is to design and implement a piece of software for playful language learning, using web technologies and the fresh CloudASR library. A 
common user interacts with their web browser almost exclusively using mouse and keyboard. Thanks to the software this thesis represents the user has an opportunity of 
delving into sometimes unpopular language learning process using his natural voice. This fact presents new and exciting possibilities, mainly regarding user 
interactivity. Big stress has been put to user friendliness, graphical fidelity and to competitive aspect of language education, exploiting Facebook integration and 
point-scoring leader boards.

Keywords:
% 3 až 5 klíčových slov v angličtině
automatic speech recognition, ASR, games, web, HTML5, Javascript

\vss}

\newpage

%%% Strana s automaticky generovaným obsahem bakalářské práce. U matematických
%%% prací je přípustné, aby seznam tabulek a zkratek, existují-li, byl umístěn
%%% na začátku práce, místo na jejím konci.

\openright
\pagestyle{plain}
\setcounter{page}{1}
\tableofcontents

%%% Jednotlivé kapitoly práce jsou pro přehlednost uloženy v samostatných souborech
\chapwithtoc{Úvod}

Základním komunikačním kanálem mezilidské interakce je mluvená řeč. Ačkoli trend v poslední době spěje spíše k separaci této komunikace virtuálními médii, je mluvené slovo i nadále nenahraditelnou součástí socializace a výměny informací. Proč tedy nepřenést schopnost porozumět lidské mluvě i na stroje, jakožto právě na ona separující virtuální média? Současná technologie je k tomu jistě dostačující, výpočetní výkon lidstva roste po zběsilé křivce a technologické know-how po ještě zběsilejší.

Aby tento cíl mohl být dosažen, je kritické mít dobrý \nom{ASR}{Automatic Speech Recognition - automatické rozpoznávání řeči} (Automatic Speech Recognition) systém. Těchto se pohybuje v éteru mnoho, volně dosažitelných, rozličné kvality a spolehlivosti. Pro specifické využití v této práci byl zvolen systém 
\verb|CloudASR|\footnote{\url{https://github.com/UFAL-DSG/cloud-asr/}}, jehož autorem je kolega Ondřej Klejch.

\texttt{\nom{TrogASR}{Translating Online Game using Automatic Speech Recognition}} (Translating Online Game using Automatic Speech Recognition) je, jak možná název napovídá, online hra, zaměřující se na překlad slov a frází z~jazyka do jazyka za využití automatizovaného rozpoznávání řeči.

Cílem této práce je návrh a implementace graficky přitažlivé hry pro webové prohlížeče, která poskytne uživatelům další, neotřelou, možnost, jak si vyzkoušet a~zdokonalit své jazykové schopnosti. Tato hra pak bude sloužit jako kompetitivitou hnaný motor pro sběr verifikovaných dat pro použitou \verb|CloudASR| knihovnu.

Velký důraz je kladen na grafickou stránku věci; uživatelsky neatraktivní hra nemá valnou naději na šíření mezi uživateli samotnými samovolně. Navíc se na~ošklivou aplikaci nikdo nebude dívat po delší časový úsek, než je jedno sezení, a~už vůbec nikdo se k takovéto aplikaci nebude vracet. Proto bylo veškeré stylování a~malování a~navrhování layoutů a přechodů prováděno s výraznou pílí.

Dále je důraz kladen i na zmíněnou kompetitivitu. Každý má touhu vidět své jméno někde vysoko na žebříčku, zanechávaje své přátele na tomto žebříčku daleko pod sebou. Proto byla zapojena i integrace s Facebookovým API, jež~přináší kýženého výsledku poměrně nenásilnou formou. Na \uv{Like} a \uv{Login} tlačítko jsou uživatelé zvyklí a~nebojí se jej extenzivně používat.

Následující text popisuje a dokumentuje vývoj této aplikace a~myšlenkové pochody při něm stojící. První kapitola analyzuje e-learning jako způsob předávání informací. Druhá kapitola zmiňuje ASR teorii, popisuje některé postupy a základní paradigmata. Třetí kapitola seznámí čtenáře s využitými technologiemi a frameworky pro vývoj aplikace \verb|trogASR|. Kapitola čtvrtá nabídne čtenáři vhled do řešení a~implementace aplikace \verb|trogASR|, do použitých návrhových vzorů a do problematiky s nimi spojené. A konečně, kapitola pátá shrne celou práci.
\chapter{Použité technologie}

\section{Kapitola}

\section{Kapitola}


\chapter{Implementace}

\section{Kapitola}

\section{Kapitola}

% Ukázka použití některých konstrukcí LateXu (odkomentujte, chcete-li)
%%% Ukázka použití některých konstrukcí LaTeXu

\subsection{Ukázka \LaTeX{}u}
\label{ssec:ukazka}

V~této krátké části ukážeme použití několika základních konstrukcí \LaTeX{}u,
které by se vám mohly při psaní práce hodit.

Třeba odrážky:

\begin{itemize}
\item Logo Matfyzu vidíme na obrázku~\ref{fig:mff}.
\item Tato subsekce má číslo~\ref{ssec:ukazka}.
\item Odkaz na literaturu
\end{itemize}

Druhy pomlček:
červeno-černý (krátká),
strana 16--22 (střední),
$45-44$ (minus),
a~toto je --- jak se asi dalo čekat --- vložená věta ohraničená dlouhými pomlčkami.
(Všimněte si, že jsme za \verb|a| napsali vlnovku místo mezery: to aby se
tam nemohl rozdělit řádek.)

% Makro na české uvozovky (novější verze LaTeXu ho už mají zabudované)
%\newcommand{\uv}[1]{\quotedblbase #1\textquotedblleft}
\uv{České uvozovky.}

\newtheorem{theorem}{Věta}
\newtheorem*{define}{Definice}	% Definice nečíslujeme, proto "*"

\begin{define}
{\sl Strom} je souvislý graf bez kružnic.
\end{define}

\begin{theorem}
Tato věta neplatí.
\end{theorem}

\begin{proof}
Neplatné věty nemají důkaz.
\end{proof}

\begin{figure}
	\centering
	\includegraphics[width=30mm]{img/logo.png}
	\caption{Logo MFF UK}
	\label{fig:mff}
\end{figure}


\chapter*{Závěr}
\addcontentsline{toc}{chapter}{Závěr}


%%% Seznam použité literatury
%%% Seznam použité literatury je zpracován podle platných standardů. Povinnou citační
%%% normou pro bakalářskou práci je ISO 690. Jména časopisů lze uvádět zkráceně, ale jen
%%% v kodifikované podobě. Všechny použité zdroje a prameny musí být řádně citovány.

\def\bibname{Seznam použité literatury}
\begin{thebibliography}{99}
\addcontentsline{toc}{chapter}{\bibname}

\bibitem{bahl_jelinek_1983}
  {\sc Bahl,} L. R., {\sc Jelinek,} F., {\sc Mercer,} R. L.\\
  \emph{A maximum likelihood approach to continuous speech recognition}.\\
  IEEE Transactions on Pattern Analysis and Machine Intelligence, 1983

\bibitem{jelinek_2010}
  {\sc Deoras,} A., {\sc Filimonov,} D., {\sc Harper,} M., {\sc Jelinek,} F.\\
  \emph{Model combination for Speech Recognition using Empirical Bayes Risk minimization}.\\
  IEEE Spoken Language Technology Workshop (SLT), 2010

\bibitem{jelinek_1976}
  {\sc Bahl,} L., {\sc Baker,} J., {\sc Cohen,} P., {\sc Dixon,} N.R., {\sc Jelinek,} F., {\sc Mercer,} R.L., {\sc Silverman,} H.F.\\
  \emph{Preliminary results on the performance of a system for the automatic recognition of continuous speech}.\\
  Acoustics, Speech, and Signal Processing, IEEE International Conference on ICASSP '76, 1976

\bibitem{mit_asr_2003}
  {\sc Glass,} James, {\sc Zue,} Victor\\
  \emph{6.345 Automatic Speech Recognition}.\\
  Massachusetts Institute of Technology: MIT OpenCourseWare, Spring 2003.\\
  License: \href{http://creativecommons.org/licenses/by-nc-sa/4.0/}{Creative Commons BY-NC-SA}

\bibitem{shaug_2003}
  {\sc O'Shaughnessy,} D.\\
  \emph{Interacting with computers by voice: automatic speech recognition and synthesis}.\\
  Proceedings of the IEEE (Volume:91, Issue: 9), 2003
  
\bibitem{shaug_2003}
  {\sc O'Shaughnessy,} D.\\
  \emph{Interacting with computers by voice: automatic speech recognition and synthesis}.\\
  Proceedings of the IEEE (Volume:91, Issue: 9), 2003
  
\bibitem{tohkura_1987}
  {\sc Tohkura,} Y.\\
  \emph{A Weighted Cepstral Distance Measure for Speech Recognition}.\\
  IEEE Trans. ASSP, Vol. ASSP-35, No. 10, 1414-1422, 1987

\bibitem{sahi_2012}
  {\sc Sahidullah,} Md., {\sc Saha,} Goutam\\
  \emph{Design, analysis and experimental evaluation of block based transformation in MFCC computation for speaker recognition}.\\
  Speech Communication, Volume 54, Issue 4, Pages 543–565, May 2012
  
\bibitem{bogert_1963}
  {\sc Bogert,} B.P., {\sc Healy,} M.J.R., {\sc Tukey,} J.W.\\
  \emph{The frequency alanysis of time series for echoes: cepstrum, pseudo-auto-covariance, cross-cepstrum and saphe cracking}.\\
  Time Series Analysis, M. Rosenblatt, Ed. New York: Wiley, ch.15, 1963

\bibitem{hossan_2010}
  {\sc Hossan,} M.A., {\sc Memon,} S., {\sc Gregory,} M.A.\\
  \emph{A novel approach for MFCC feature extraction}.\\
  Signal Processing and Communication Systems (ICSPCS), 4th International Conference, 2010
  
\bibitem{byeo_2012}
  {\sc Byeongkyu,} Ko, {\sc Dongjin,} Choi, {\sc Chang,} Choi, {\sc Junho,} Choi, {\sc Pankoo,} Kim\\
  \emph{Document Classification through Building Specified N-Gram}.\\
  Innovative Mobile and Internet Services in Ubiquitous Computing (IMIS), Sixth International Conference 2012

\bibitem{masa_1997}
  {\sc Masataki,} H., {\sc Sagisaka,} Y., {\sc Hisaki,} K., {\sc Kawahara,} T.\\
  \emph{Task adaptation using MAP estimation in N-gram language modeling}.\\
  Acoustics, Speech, and Signal Processing (vol.2), IEEE International Conference 1997

\bibitem{jelinek_1980}
  {\sc Jelinek,} F., {\sc Mercer,} R.\\
  \emph{Interpolated estimation of Markov source parameters from sparse data}.\\
  Proceedings of Workshop on Pattern Recognition in Practice, pg.381-397, 1980

\bibitem{katz_1987}
  {\sc Katz,} S.\\
  \emph{Estimation of probabilities from sparse data for the language model component of a speech recognizer}.\\
  Acoustics, Speech and Signal Processing, IEEE Transactions, 1987

\bibitem{church_1991}
  {\sc Church,} K., {\sc Gale,} W.\\
  \emph{A Comparison of the Enhanced Good-Turing and Deleted Estimation Methods for Estimating Probabilities of English Bigrams}.\\
  Computer Speech \& Language, 1991

\bibitem{chienli_2013}
  {\sc Chien-Lin,} Huang, {\sc Hori,} C., {\sc Kashioka,} H.\\
  \emph{Semantic inference based on neural probabilistic language modeling for speech indexing}.\\
  Acoustics, Speech and Signal Processing (ICASSP), IEEE International Conference 2013

\bibitem{goblirsch_1996}
  {\sc Goblirsch,} D.M.\\
  \emph{Viterbi beam search with layered bigrams}.\\
  Spoken Language, Fourth International Conference, 1996

\bibitem{poritz_1988}
  {\sc Poritz,} A.\\
  \emph{Hidden Markov models: a guided tour}.\\
  Acoustics, Speech, and Signal Processing, International Conference, 1988

\bibitem{han_2003}
  {\sc Han,} Shu, {\sc Hetherington,} I.L., {\sc Glass,} J.\\
  \emph{Baum-Welch training for segment-based speech recognition}.\\
  Automatic Speech Recognition and Understanding, IEEE Workshop 2003

\bibitem{galindo_1995}
  {\sc Galindo,} P.L.\\
  \emph{The competitive forward-backward algorithm (CFB)}.\\
  Artificial Neural Networks, Fourth International Conference, 1995
  
\bibitem{kingsbury_2009}
  {\sc Kingsbury,} B.\\
  \emph{Lattice-based optimization of sequence classification criteria for neural-network acoustic modeling}.\\
  Acoustics, Speech and Signal Processing, IEEE International Conference 2009

\bibitem{tarar_2010}
  {\sc Tarar,} S.\\
  \emph{Speech analysis: Desktop items activation using Dynamic time warping algorithm}.\\
  Computer Science and Information Technology (ICCSIT), 3rd IEEE International Conference 2010 

\bibitem{franzini_1990}
  {\sc Franzini,} Michael, {\sc Lee,} K-F, {\sc Waibel,} Alex\\
  \emph{Connectionist viterby training: A new hybrid method for continuous speech recognition}.\\
  Acoustics, Speech, and Signal Processing, 1990.ICASSP-90., pp.452-428, International Conference, 1990
  
\bibitem{rasipuram_2014}
  {\sc Rasipuram,} Ramya, {\sc Magimai.-Doss,} Mathew, \\
  \emph{Acoustic and Lexical Resource Constrained ASR using Language-Independent Acoustic Model and Language-Dependent Probabilistic Lexical Model}.\\
  Idiap Research Report, 02-2014

\bibitem{haili_2013}
  {\sc Haili,} Luo\\
  \emph{The research of applying regular grammar to making model for lexical analyzer}.\\
  Information Management, Innovation Management and Industrial Engineering (ICIII), 6th International Conference, 2013
  
\bibitem{schluter_2010}
  {\sc Schluter,} R., {\sc Nussbaum-Thom,} M., {\sc Ney,} H.\\
  \emph{On the Relationship Between Bayes Risk and Word Error Rate in ASR}.\\
  Audio, Speech, and Language Processing, IEEE Transactions, 2010

\bibitem{korviny_2005}
  {\sc Korviny,} Petr\\
  \emph{Moodle (nejen) na OPF}.\\
  OPF, 2005

\bibitem{stockley_2003}
  {\sc Stockley,} Derek\\
  http://www.derekstockley.com.au/elearning-definition.html, 2003

\bibitem{prucha_2009}
  {\sc Průcha,} Jan, {\sc Walterová,} Eliška, {\sc Mareš,} Jiří\\
  \emph{Pedagogický slovník}.\\
  Pedagogický slovník s.66, Praha, Portál, 2003

\bibitem{fongling_2006}
  {\sc Fongling,} Fu, {\sc Sheng-Chin,} Yu\\
  \emph{The Games in E Learning Improve the Performance}.\\
  Information Technology Based Higher Education and Training, 7th International Conference, 2006

\bibitem{graven_2006}
  {\sc Graven,} O.H., {\sc MacKinnon,} L.\\
  \emph{Exploitation of games and virtual environments for e-learning}.\\
  Information Technology Based Higher Education and Training, 7th International Conference, 2006

\bibitem{xuemin_2009}
  {\sc Xue-min,} Zhang, {\sc Mao,} Li, {\sc Bin,} Yang, {\sc Liu,} Chang\\
  \emph{Violent Components and Interactive Mode of Computer Video Game on Player's Negative Social Effect}.\\
  Intelligent Information Technology Application, Third International Symposium, 2009

\bibitem{bialystok_2006}
  {\sc Bialystok,} Ellen\\
  \emph{Effect of bilingualism and computer video game experience on the Simon task}.\\
  Canadian Journal of Experimental Psychology/Revue canadienne de psychologie expérimentale, Vol 60(1), 68-79, Mar 2006
  
\bibitem{ahn_2006}
  {\sc von Ahn,} Luis\\
  \emph{Games With A Purpose}.\\
  Institute of Electrical and Electronics Engineers Computer Magazine, vol. 39, no. 6, 2006

\bibitem{ahn_2004}
  {\sc von Ahn,} Luis, {\sc Dabbish,} Laura\\
  \emph{Labeling images with a computer game}.\\
  In Proceedings of the SIGCHI Conference on Human Factors in Computing Systems, pp. 319-326, 2004
 
\bibitem{ahn_2008}
  {\sc von Ahn,} Luis, {\sc Maurer,} Benjamin, {\sc McMillen,} Colin, {\sc Abraham,} David, {\sc Blum,} Manuel\\
  \emph{reCAPTCHA: Human-Based Character Recognition via Web Security Measures}.\\
  Science, vol. 321, no 5895, pp. 1465-1468, 2008
 
\end{thebibliography}


%%% Tabulky v bakalářské práci, existují-li.
\chapwithtoc{Seznam tabulek}

%%% Použité zkratky v bakalářské práci, existují-li, včetně jejich vysvětlení.
\chapwithtoc{Seznam použitých zkratek}
\textbf{ASR} -- Automatic Speech Recognition -- automatické rozpoznávání řeči

%%% Přílohy k bakalářské práci, existují-li (různé dodatky jako výpisy programů,
%%% diagramy apod.). Každá příloha musí být alespoň jednou odkazována z vlastního
%%% textu práce. Přílohy se číslují.
\chapwithtoc{Přílohy}

\openright
\end{document}
